% Run GCT_KRON_20261120; GPT-6 Astra Pro; effort not exposed.
\section{Actual positive highest-weight algebras on the Hessian ray}

\subsection*{A-034. Every stretch gives a strict actual multiplicity gap}
For \(n\geq3\), retain the notation
\[
 R=n^2,\qquad D=R+1,\qquad q=R(n-2),\qquad
 E=n+q=n(n-1)^2,\qquad
 \lambda_n=(E+1,3,2^{R-2}).
\]
Define the positive partition functions
\[
 p_n(s)=[z^s]\prod_{i=1}^n(1-z^i)^{-1},
 \qquad
 h_n(s)=[z^s]\prod_{i=2}^n(1-z^i)^{-1}.
                                                               \tag{A-034.1}
\]
Then, for every \(s\geq1\),
\[
 \boxed{\begin{aligned}
 q^n_{s\lambda_n}(sD[n])&\geq p_n(s)>h_n(s),\\
 K_n(s\lambda_n)
 &=\operatorname{sk}\bigl(s\lambda_n,((sD)^n)\bigr)\\
 &=g\bigl(s\lambda_n,((sD)^n),((sD)^n)\bigr)=h_n(s).
 \end{aligned}}                                                \tag{A-034.2}
\]
In particular the actual determinant and permanent
multiplicities are both positive for every \(s\geq2\), and
\[
 q^n_{s\lambda_n}(sD[n])-K_n(s\lambda_n)
                         \geq p_n(s-1)>0.                       \tag{A-034.3}
\]
At \(s=1\) this recovers the occurrence obstruction of A-030:
the determinant multiplicity is zero and the permanent
multiplicity is one. The proof of (A-034.2) constructs actual
polynomial highest-weight functions, uses no finite rank
certificate, and works for all \(n\geq3\).
It concerns the unpadded case \(n=m\).

There is a stronger algebraic statement. Write
\(\mathcal D_n=\overline{\operatorname{GL}_R\cdot\det_n}\),
\(\mathcal P_n=\overline{\operatorname{GL}_R\cdot\operatorname{per}_n}\),
and let \(U\) be the upper unitriangular subgroup with the
Borel convention of A-033. For \(\mathcal X=\mathcal D_n\)
or \(\mathcal P_n\), put
\[
 \mathcal R_\lambda(\mathcal X)
 =\bigoplus_{s\geq0}
       \mathbb C[\mathcal X]_{sD}^{\,U}[s\lambda_n].
                                                               \tag{A-034.4}
\]
The square brackets specify the torus weight space.
There are explicit universal polynomials \(Z,U_2,\ldots,U_n\)
of ray degrees \(1,2,\ldots,n\) such that
\[
 \boxed{\begin{gathered}
 \mathbb C[Z,U_2,\ldots,U_n]
       \lhook\joinrel\longrightarrow \mathcal R_\lambda(\mathcal P_n),
 \\[2pt]
 \mathcal R_\lambda(\mathcal D_n)
                         =\mathbb C[U_2,\ldots,U_n].
 \end{gathered}}                                                \tag{A-034.5}
\]
Both displayed polynomial algebras have algebraically
independent named generators. Restriction of the universal
subalgebra to the determinant closure has kernel \((Z)\).
This describes the two restrictions of that universal
subalgebra; it does not assert a restriction homomorphism
between the entire permanent and determinant coordinate rings.

\paragraph{Normalized binary coefficients and the generators.}
For a degree-\(n\) form \(F\) in \(R\) variables, let
\(H(F)=\det\operatorname{Hess}_{R}(F)\), the Hessian in all
variables. Define \(F_i,H_i\) by
\[
 \begin{split}
 F(xe_1+ye_2)&=\sum_{i=0}^n\binom ni F_i x^{n-i}y^i,\\
 H(F)(xe_1+ye_2)&=\sum_{i=0}^q\binom qi H_i x^{q-i}y^i.
 \end{split}
                                                               \tag{A-034.6}
\]
Thus \(F_i\) has coefficient degree one and \(H_i\) has
coefficient degree \(R\). Set
\[
 \begin{split}
 I_i={}&(-1)^{i-1}(i-1)F_1^i\\
       &+\sum_{j=2}^i(-1)^{i-j}\binom ij
                    F_0^{j-1}F_1^{i-j}F_j
                    \qquad(2\leq i\leq n),\\
 Z={}&F_0H_1-F_1H_0,\qquad
 U_i=H_0^i I_i\qquad(2\leq i\leq n).
 \end{split}                                                   \tag{A-034.7}
\]
These are polynomials, with no localization in their
definition. The centered expression
\[
 I_i=F_0^{i-1}
       \sum_{j=0}^i\binom ij(-F_1/F_0)^{i-j}F_j
                                                               \tag{A-034.8}
\]
is an identity on \(F_0\ne0\); its \(j=0,1\) terms combine
to the first polynomial term in (A-034.7).
It proves that \(I_i\) is invariant under the binary
unipotent substitution \(x\mapsto x+t y\).
Indeed \(F_1/F_0\) changes to \(F_1/F_0+t\), exactly
canceling that substitution. The resulting invariance is
a polynomial identity and holds also on \(F_0=0\).

For an upper triangular \(g\), write
\(ge_1=a e_1\), \(ge_2=b e_1+c e_2\).
The normalized coefficient rule in A-033 and
\(H(F\circ g)=\det(g)^2H(F)\circ g\) give
\[
 \begin{split}
 I_i(F\circ g)&=a^{i(n-1)}c^i I_i(F),\\
 H_0(F\circ g)&=\det(g)^2a^q H_0(F),\\
 Z(F\circ g)&=\det(g)^2a^{E-1}c\,Z(F).
 \end{split}                                                   \tag{A-034.9}
\]
For \(I_i\), the diagonal character can also be checked
term by term in (A-034.7).
Consequently \(Z\) has coefficient degree \(D\) and
highest weight \(\lambda_n\); \(U_i\) has coefficient degree
\(iD\) and highest weight \(i\lambda_n\).
The polynomial \(Z\) is a nonzero scalar multiple of the
A-030 Hessian-bracket highest-weight function:
for any fixed form \(P\) and pair \(A,B\),
\[
 Z=\frac1q P(A)\,dH(P)_A(B)
                   -\frac1n H(P)(A)\,dP_A(B).                 \tag{A-034.10}
\]
Here and below binary coefficients at \(A,B\) mean those
of the pencil \(xA+yB\).

\paragraph{Binary-pencil evaluations are genuine orbit evaluations.}
If \(A,B\in W=\operatorname{Mat}_n(\mathbb C)\) are linearly
independent, extend them to the first two columns of an
element \(g\in\operatorname{SL}(W)\).
Such an extension exists: complete a basis arbitrarily and
scale one remaining column to make its determinant one.
Thus the coefficients of
\[
 P(xA+yB),\qquad H(P)(xA+yB)
\]
are exactly the functions in (A-034.6) evaluated at
\(P\circ g\), with no additional Hessian factor.
This observation applies to \(P=\operatorname{per}_n\) and
to \(P=\det_n\).
A polynomial relation on either orbit closure therefore
pulls back to a relation on the open set of independent
pairs, and hence to an identity in the polynomial ring
\(\mathbb C[A,B]\). No chosen basis completion must vary
regularly with \(A,B\) for this argument.

\paragraph{A fixed permanent Hessian zero with dominant binary coefficients.}
Let \(P=\operatorname{per}_n\) and let \(A_*\) be the
all-ones matrix except for its \(11\) entry \(3-n\).
The permanent sum and the exact Hessian identity (A-030.13)
give
\[
        P(A_*)=2(n-1)!\ne0,\qquad H(P)(A_*)=0.                 \tag{A-034.11}
\]
The same Hessian identity at an \(11\) entry other than
\(3-n\) shows that \(H(P)\) is not the zero polynomial.

For \(B=\operatorname{diag}(b_1,\ldots,b_n)\), write
\[
 P(A_*+tB)=\sum_{k=0}^n c_k t^k,
 \qquad e_j=e_j(b_2,\ldots,b_n),\quad e_0=1.
\]
Selecting the \(k\) diagonal entries contributed by \(tB\)
in the defining permanent sum gives
\[
 \begin{split}
 c_0&=2(n-1)!,\\
 c_k&=(n-k)!\,b_1 e_{k-1}
            +(2-k)(n-k-1)!\,e_k\quad(1\leq k\leq n-1),\\
 c_n&=b_1e_{n-1}.
 \end{split}                                                   \tag{A-034.12}
\]
If the selected set contains index \(1\), the remaining
matrix is all ones. Otherwise its remaining exceptional
entry is still \(3-n\), and its permanent is
\((2-k)(n-k-1)!\). This proves every coefficient in
(A-034.12), including the vanishing second summand for
\(k=2\).

The first two variable coefficients are
\[
 c_1=(n-1)!\,b_1+(n-2)!\,e_1,\qquad
 c_2=(n-2)!\,b_1e_1.
 \]
In the variables \(b_1,e_1,\ldots,e_{n-1}\), the Jacobian
of \((c_1,\ldots,c_n)\) has its first two rows supported
in the first two columns. The remaining block, with rows
\(c_3,\ldots,c_n\) and columns \(e_2,\ldots,e_{n-1}\),
is upper triangular, with diagonal entries
\((n-k)!\,b_1\), for \(3\leq k\leq n\).
Its full determinant is therefore
\[
 ((n-2)!)^2\bigl((n-1)b_1-e_1\bigr)
              b_1^{n-2}\prod_{k=3}^n(n-k)! \ne0
                          \quad\hbox{as a polynomial}.        \tag{A-034.13}
\]
The elementary-symmetric map
\((b_2,\ldots,b_n)\mapsto(e_1,\ldots,e_{n-1})\)
is dominant over \(\mathbb C\), since every monic
degree-\((n-1)\) polynomial has all its roots there.
It follows from (A-034.13) that the variable coefficients
\((c_1,\ldots,c_n)\) range over a dense subset of
\(\mathbb A^n\), with \(c_0\) fixed.

On this family \(F_0=2(n-1)!\), and
\(F_i=c_i/\binom ni\). The triangular polynomial map
\[
     (F_1,F_2,\ldots,F_n)
               \longmapsto(F_1,I_2,\ldots,I_n)
 \]
has nonzero diagonal coefficients
\(1,F_0,F_0^2,\ldots,F_0^{n-1}\).
Consequently \(I_2,\ldots,I_n\), evaluated at
\((A_*,B)\), are algebraically independent.
In particular, for every nonzero polynomial \(T\),
\[
                 T(I_2,\ldots,I_n)(A_*,B)
                      \quad\hbox{is not identically zero}.    \tag{A-034.14}
\]
Since \(A_*\) has nonzero off-diagonal entries, it is
independent of every nonzero diagonal \(B\); the dense
families used here consist of valid orbit frames.

\paragraph{A prime valuation separates every permanent monomial.}
Factor the nonzero polynomial \(H(P)(A)\) in
\(\mathbb C[A]\). By (A-034.11), some irreducible factor
\(h(A)\) vanishes at \(A_*\). Write
\[
                 H(P)(A)=h(A)^e v(A),\qquad
                 e\geq1,\qquad h\nmid v.
\]
Because \(P(A_*)\ne0\), we have \(h\nmid P\).
Moreover (A-034.14) implies
\[
       h\nmid T(I_2,\ldots,I_n)\quad\hbox{in }\mathbb C[A,B]
                         \quad\hbox{whenever }T\ne0.           \tag{A-034.15}
\]
Indeed divisibility would make its entire specialization
at \(A_*\) vanish, contrary to (A-034.14).

Let \(\operatorname{ord}_h\) be the prime-factor valuation
on the fraction field of \(\mathbb C[A,B]\).
In characteristic zero, some partial derivative of the
nonconstant irreducible polynomial \(h\) is nonzero and
is not divisible by \(h\), by its smaller degree.
Thus \(d h_A(B)\) is nonzero modulo \(h\).
Differentiating \(H(P)=h^e v\) gives
\[
 dH(P)_A(B)=e h^{e-1}v\,dh_A(B)+h^e\,dv_A(B).
\]
In (A-034.10), the first term has valuation \(e-1\),
while the second has valuation at least \(e\).
Since \(h\nmid P\), it follows that
\[
                         \operatorname{ord}_h Z=e-1.          \tag{A-034.16}
\]
This is a valuation at a prime divisor, not a claim that
the gradient of \(H(P)\) is nonzero at the particular
point \(A_*\).

Suppose there were a nonzero polynomial relation among
\(Z,U_2,\ldots,U_n\) in the permanent coordinate ring.
The coefficient-degree grading allows one to take a
nonzero weighted-homogeneous relation of degree \(s\),
where the generator weights are \(1,2,\ldots,n\).
Write it as
\[
                    \sum_j Z^j T_j(U_2,\ldots,U_n)=0,
\]
with each nonzero \(T_j\) weighted homogeneous of
degree \(s-j\). On binary pencils this becomes
\[
             \sum_j Z^j H(P)(A)^{s-j}
                                T_j(I_2,\ldots,I_n)=0.         \tag{A-034.17}
\]
By (A-034.15)--(A-034.16), the valuation of its
nonzero \(j\)-th summand is
\[
                  j(e-1)+e(s-j)=es-j.
\]
These valuations are distinct. The summand with largest
\(j\) has a uniquely smallest valuation and cannot cancel
in a zero sum. This contradicts (A-034.17).
It proves the permanent injection in (A-034.5).
Counting its monomials of weighted degree \(s\) gives
the lower bound \(q^n_{s\lambda_n}(sD[n])\geq p_n(s)\).

\paragraph{The determinant ray algebra is exactly polynomial.}
The same \(I_2,\ldots,I_n\) are algebraically independent
on determinant binary pencils. Indeed diagonal pencils
realize every product of \(n\) binary linear forms, and
every binary degree-\(n\) form factors into such a product
over \(\mathbb C\). One can also keep \(F_0=1\) by
using \(A=I_n\), \(B=\operatorname{diag}(b_1,\ldots,b_n)\):
the coefficients are elementary symmetric polynomials,
and (A-034.7) is again triangular with nonzero diagonal.
Restricting to the dense open set where \(A,B\) are
independent makes them genuine special-linear orbit
evaluations.

The full Hessian identity in A-030 is
\[
 H(\det_n)=c_n(\det_n)^{n(n-2)},\qquad c_n\ne0.
                                                               \tag{A-034.18}
\]
In particular \(H_0\) is a nonzero polynomial function
on \(\mathcal D_n\), since its value on a frame with
first column \(I_n\) is \(c_n\).
The monomials in \(I_2,\ldots,I_n\) of weighted degree
\(s\) form an independent space of dimension \(h_n(s)\).
Their coefficient degree is \(s\) and their highest
weight is \((s(n-1),s)\).
The ring \(\mathbb C[\mathcal D_n]\) is a domain, so
multiplication of that entire space by \(H_0^s\) is
injective. The resulting monomials in \(U_2,\ldots,U_n\)
have degree \(sD\), highest weight \(s\lambda_n\), and
are independent. Hence
\[
                         K_n(s\lambda_n)\geq h_n(s).
                                                               \tag{A-034.19}
\]
The exact ordinary Kronecker formula in A-030.15,
obtained by the determinant twist to the two-row core,
is \(g(s\lambda_n,((sD)^n),((sD)^n))=h_n(s)\).
Combining this with \(K_n\leq\operatorname{sk}\leq g\)
gives equality throughout (A-034.2).
At every degree the independent monomials in
\(U_2,\ldots,U_n\) already have this full dimension,
so they span the entire ray highest-weight space.
This proves the determinant equality in (A-034.5),
including algebraic independence.

Finally (A-034.18), with \(q=n^2(n-2)\), gives
\(H_1=c_n F_0^{n(n-2)-1}F_1\) on raw determinant
pencils. For \(F=\det_n\circ g\) on the full
general-linear orbit, Hessian covariance multiplies
both \(H_0\) and \(H_1\) by the same additional
factor \(\det(g)^2\). Hence the two terms of \(Z\)
still cancel there, and \(Z=0\) on the determinant
closure by polynomiality. Independence of the \(U_i\)
shows that the kernel on the universal polynomial
subalgebra is precisely \((Z)\), as asserted.

\paragraph{Strictness, positive support, and relative saturation.}
The generating functions in (A-034.1) satisfy
\[
             h_n(s)=p_n(s)-p_n(s-1)\qquad(s\geq1).
 \]
Since \(p_n(s-1)\geq1\), this proves (A-034.3).
The parts \(2,3\) are available for every \(n\geq3\);
therefore \(h_n(1)=0\) and \(h_n(s)>0\) for every
\(s\geq2\). Both sides of the actual multiplicity
obstruction are consequently positive at every such
stretch, with no unspecified eventual threshold.
The three-copy bound in A-033 remains stronger than
the present lower bound at \(s=2\).

The exact determinant ray algebra is a polynomial
algebra with positive monomial basis and generator
degrees \(2,3,\ldots,n\). Its degree support is
\(\{0,2,3,4,\ldots\}\); its generated degree group
is \(\mathbb Z\), because \(\gcd(2,3)=1\).
Thus degree \(1\) is a relative-saturation hole
despite the algebra being normal and having an
explicit positive monomial basis. Equivalently,
\(2\lambda_n\) and \(3\lambda_n\) occur on the
determinant closure, so \(\lambda_n\) belongs to
the group generated by its actual support; it
does not itself occur. This realizes the positive
weighted-polynomial model of O-003 inside an actual
determinant highest-weight algebra.
Normality of an algebra with multidimensional weight
spaces does not make its weight semigroup saturated;
the multiplicity-free hypothesis in A-019 cannot be
discarded.

\paragraph{Geometric dictionary, sources, and scope.}
All generators above are actual polynomial functions
over \(\mathbb Q\), after the displayed nonzero
factorial denominators are allowed. The Hessian
covariance, centered coefficients, and multiplication
are polynomial comodule operations in the sense of
SGA 3, I, Proposition 4.7.2, pp.~34--35; their
highest-weight characters use Proposition 4.7.3,
pp.~35--36. The Schur and exterior operations are
those of SGA 6, VI, Theorem 3.3, pp.~392--393;
the characteristic-zero irreducible interpretation
uses Macdonald, I, Appendix A(5.3)--(5.6),
pp.~155--156, as in F-016.
The known determinant Hessian identity, together
with its normalization and the explicit permanent
Hessian formula used at \(A_*\), is documented and
proved in A-030. The ordinary two-row Kronecker
calculation is A-030.15, using O-002 and F-002.
The standard orbit bound \(K_n\leq\operatorname{sk}\leq g\)
is the already cited determinant stabilizer formula.
No historical novelty claim is made.

The exact determinant conclusion here strengthens
the ray consequences of O-028: no invariant degree
monoid or numerical invariant nonvanishing is needed
for this particular family. Its broader invariant
translation statements remain applicable elsewhere.
For \(s>0\), the partition \(s\lambda_n\) has length
\(n^2\), whereas the internally padded permanent
with \(m<n\) uses at most \(m^2\) variables.
Its multiplicity at this weight is therefore zero.
These results give an actual family at \(n=m\);
they imply no large-padding lower bound or resolution
of permanent-versus-determinant complexity.
